% ====================================================================
% §17 MSMR 동형과 전자항 plug-in — LCO 로의 일반화
% ====================================================================
\section{MSMR 동형과 전자항 plug-in --- LCO 로의 일반화 (파라미터 교체 $+$ 전자항 추가)}\label{sec:lco-code}
이 통합의 실용적 결론 --- ``같은 골격이 LCO 곡선을 낸다'' --- 의 근거는 \emph{MSMR 동형}이다:
MSMR 모델이 양극 \emph{조성}을 전위의 전이별 logistic 합으로 적기 때문이다. 계보를 한 줄로 짚으면 ---
삽입 전극의 개방회로 전위를 종(species)별 Nernst-형 점유 합으로 닫는 치환 격자 열역학 모델을 Verbrugge
등이 흑연$\cdot$스피넬$\cdot$인산철$\cdot$층상 산화물에 한 틀로 세웠고\cite{msmr_origin2017}, Baker--Verbrugge
의 다공 전극 정식화가 이를 multi-species, multi-reaction(MSMR) 모델로 명명했으며\cite{bakerverbrugge2018},
엔트로피$\cdot$온도 의존 파라미터화가 뒤따랐다\cite{msmr2024}.

\subsection{MSMR 대응 --- 같은 logistic, 같은 부호 규약}\label{sec:lco-code-msmr}
\textbf{(a) 출발 --- MSMR 종별 점유식.}
\begin{equation}
x_j=\frac{X_j}{1+\exp[\,f(U-U_j^0)/\omega_j\,]},
\label{eq:msmr}
\end{equation}
여기서 $X_j$ 는 종(전이) $j$ 의 용량 분율, $U_j^0$ 는 종 중심 전위다. 원계열 MSMR 의 지수 인자
$f=F/RT$(항상 양)$\cdot$무차원 폭 $\omega_j$ 를 본 장은 $F/RT$ 를 폭에 흡수($\omega_j\mapsto\omega_jRT/F$,
전압 차원)해 방향 부호 $f=\pm1$ 만 지수에 남기도록 재모수화하며, 이는 폭 슬롯 대응이 흑연
식~\eqref{eq:wbase} 의 $w_j=n_jRT/F$($n_j\leftrightarrow$무차원 $\omega_j$)와 동형임을 드러낸다(정식
대응은 아래 식~\eqref{eq:lco-msmrmap}).
\begin{srcbox}
\textbf{다리 --- eq:msmr 가 어느 정식화이고 계보의 어디서 명명됐는가.} 식~\eqref{eq:msmr} 의 종별
logistic 점유식은 Verbrugge 등\cite{msmr_origin2017} 이 삽입 전극 OCV 를 치환 격자 열역학의 종별
Nernst-형 점유 합으로 닫아 세운 표준 파라미터화이고, Baker--Verbrugge\cite{bakerverbrugge2018} 의 다공
전극 정식화가 이를 ``multi-species, multi-reaction(MSMR)'' 로 명명한 것이다. 대응은 다음과 같다(왼쪽이
본문 기호, 오른쪽이 원 논문 량):
\begin{center}\footnotesize
\begin{tabular}{@{}ll@{}}
\hline
본문(식~\eqref{eq:msmr},~\eqref{eq:lco-msmrmap}) & MSMR\cite{msmr_origin2017,bakerverbrugge2018} 대응량 \\
\hline
종별 점유 $x_j=X_j/(1+e^{f(U-U_j^0)/\omega_j})$ & 동일 종별 점유식(Verbrugge 2017) \\
지수 인자 $f=F/RT$ ($>0$) & $f=F/RT$ (동일 정의) \\
종 용량 분율 $X_j$ & $X_j$ = 종 $j$ 가 점유 가능한 host 자리 분율 \\
종 중심 전위 $U_j^0$ & $U_j^0$ = 농도무관 표준 전극 전위 \\
무차원 폭 $\omega_j$ & $\omega_j$ = 반응 $j$ 의 무질서 정도 파라미터 \\
총 조성 $\sum_jX_j\theta_j^{\mathrm{MSMR}}$ & $x=\sum_j x_j$ (자리 보존 총합) \\
명명 ``MSMR'' & Baker--Verbrugge 2018 ``multi-species, multi-reaction'' \\
\hline
\end{tabular}
\end{center}
등치의 핵은 폭 슬롯의 재모수화다 --- 원계열의 무차원 폭 $\omega_j$ 에서 $F/RT$ 를 폭으로 흡수하면
본문의 전압-차원 폭 $w_j=n_jRT/F$ 가 되고, 지수엔 방향 부호만 남는다:
\begin{equation}
\frac{f\,(U-U_j^0)}{\omega_j}\ \xrightarrow[\ \omega_j\mapsto w_j=n_jRT/F\ ]{\ f=F/RT\ \text{흡수}\ }\
\frac{\sigma_d\,(V-U_j^{\,d})}{w_j},
\qquad (U,U_j^0,\omega_j,X_j,f)\ \leftrightarrow\ (V,U_j^{\,d},w_j,Q_j,\sigma_d),
\label{eq:br-msmr-1}
\end{equation}
곧 식~\eqref{eq:lco-msmrmap} 의 1:1 대응이며, 계수비교 한 줄($f/\omega_j=\sigma_d/w_j$)이 방향 인자
$f=+\sigma_d$ 를 유일하게 고정한다(본문 §17.1(c)). \emph{방법 요지} --- Verbrugge 등은 삽입 전극 OCV 를
치환 격자 열역학의 종별 점유(자리 보존 $\sum_jX_j$)의 Nernst-형 합으로 닫아 식~\eqref{eq:msmr} 을
세웠고, Baker--Verbrugge 가 이를 다공 전극 다반응 정식화로 확장하며 MSMR 로 명명했다. \emph{가정 차}
--- 원계열 MSMR 은 방향 없는 \emph{평형 등온선}이고 폭이 무차원 $\omega_j$(지수의 $F/RT$ 항상 양)이나,
본문은 $F/RT$ 를 폭에 흡수($\omega_j\!\mapsto\!w_j$, 전압 차원)하고 방향 부호 $f=+\sigma_d$ 를 남겨
충$\cdot$방전을 재배선하며, 순정 MSMR 에 없는 히스테리시스 분기($\gamma_j$, \S\ref{sec:lco-hys})와 T1
전자 엔트로피 항(\S\ref{sec:lco-electronic})을 추가로 얹는다 --- \emph{함수형 동형이지 물리량 동일이
아니다}($x_j/X_j$ 는 리튬화 점유, $\xi$ 는 탈리튬화 진행률).
\end{srcbox}

\textbf{(b) 연산 --- 정규화.} 식~\eqref{eq:msmr} 을 종 용량으로 나누면
종별 \emph{점유율}(리튬화 분율)과 그 여집합(탈리튬화 진행률)이 나온다:
\begin{equation}
\theta_j^{\mathrm{MSMR}}\equiv\frac{x_j}{X_j}=\frac{1}{1+e^{+f(U-U_j^0)/\omega_j}},\qquad
\xi_j^{\mathrm{MSMR}}\equiv1-\frac{x_j}{X_j}=\frac{1}{1+e^{-f(U-U_j^0)/\omega_j}},
\label{eq:lco-msmrnorm}
\end{equation}
이고 총 조성은 $\sum_jX_j\theta_j^{\mathrm{MSMR}}$(리튬 함량), 추출 전하는 그 여집합의 가중합
$\sum_jQ_j\xi_j$ --- 용량 가중 $X_j\!\leftrightarrow\!Q_j$ 로 읽으면 \S\ref{sec:lco-peak} 의 전하
보존식~\eqref{eq:lco-charge} 과 같은 구조(용량-가중 logistic 합)다.
\textbf{(c) 중간식 --- 여집합 항등식과 방향 인자 대응.} Ch1 logistic 서식은 여집합에 대해 닫혀 있다 ---
같은 중심$\cdot$폭에서
\begin{equation}
1-\xi_{\eq,j}^{(\sigma_d)}(V)
=\frac{1}{1+e^{+\sigma_d(V-U_j^{\,d})/w_j}}
=\xi_{\eq,j}^{(-\sigma_d)}(V),
\label{eq:lco-comp}
\end{equation}
곧 점유와 진행률은 방향 인자 부호 하나로 오가며 종 $\xi(1-\xi)$ 는 이 교환에 불변이다(\S\ref{sec:lco-peak}(b) 의 종 항등식 도출부). 이제 \emph{같은 물리량끼리} 맞댄다 --- MSMR 점유율
$\theta_j^{\mathrm{MSMR}}=x_j/X_j$(리튬화 분율)는 Ch1 의 점유 $\theta_{\eq,j}=1-\xi_{\eq,j}$ 와, 그 여집합
$\xi_j^{\mathrm{MSMR}}$(탈리튬화 진행률)는 Ch1 의 진행률 $\xi_{\eq,j}$(식~\eqref{eq:xieq})와 같은 종류의
양이다. 진행률끼리 지수를 맞대면($-f\leftrightarrow-\sigma_d$), 점유끼리 맞대도
($+f\leftrightarrow+\sigma_d$) 동일하게, 슬롯별로
\begin{equation}
U\leftrightarrow V,\qquad U_j^0\leftrightarrow U_j^{\,d},\qquad
\omega_j\leftrightarrow w_j,\qquad X_j\leftrightarrow Q_j,\qquad
f=+\sigma_d
\label{eq:lco-msmrmap}
\end{equation}
가 1:1 로 읽힌다 --- 방향 인자 대응 $f=+\sigma_d$ 는 두 지수가 \emph{모든} $U$ 에서 일치할 것을 요구하는
항등식의, 이 pairing 하에서의 유일해다(계수비교 한 줄: 두 지수 $-f(U-U_j^0)/\omega_j$ 와
$-\sigma_d(V-U_j^{\,d})/w_j$ 가 모든 $U$ 에서 같으려면 $U$ 의 1차 계수가 같아야 하므로
$f/\omega_j=\sigma_d/w_j$ 이고, 폭 슬롯 대응 $\omega_j\leftrightarrow w_j$ 아래에서 $f=+\sigma_d$ 만 남는다). 원계열의 $f=F/RT>0$ 는 이 규약에서 탈리튬화 진행($\sigma_d=+1$ 슬롯)에 대응하며,
MSMR 평형 등온선 자체는 방향이 없는 양이라 $\sigma_d$ 는 훑는 방향의 해석만 정한다. \textbf{★같은 logistic 이라는 것은 \emph{함수형 동형이지 물리량 동일이
아니다}} --- $x_j/X_j$ 는 리튬 함량 몫(점유)이고 $\xi$ 는 탈리튬화 진행률이라, 둘을 직접 등치하면 부호가
뒤집힌 $f=-\sigma_d$ 가 나온다(식~\eqref{eq:lco-comp} 의 여집합 뒤바꿈 --- 종 $\xi(1-\xi)$ 는 이 교환에 불변이라 관측 peak 은 어느
쪽으로도 같으나, 방향 서술은 그렇지 않다). 이로써 탈리튬화/리튬화 진행 방향이 두 모델에서 일관되게
맞춰진다($U_j^0\leftrightarrow U_j^{\,d}$ 대응은 방향분해 파라미터화로 읽으며, 순정 MSMR 평형 등온선은
$\gamma_j{=}0$ 의 $U_j$ 대응이 기본형이다). \textbf{(d) 박스 --- MSMR$\to$평형 peak 변환 폐쇄.}
대응표~\eqref{eq:lco-msmrmap} 아래에서 MSMR 종별 미분용량과 Ch1 평형 peak 은 같은 식이다:
\begin{equation}
\boxed{\;Q_j\Big|\frac{\dd(x_j/X_j)}{\dd U}\Big|
=Q_j\,\frac{\theta_j^{\mathrm{MSMR}}(1-\theta_j^{\mathrm{MSMR}})}{\omega_j}
\;\equiv\;Q_j\,\frac{\xi_{\eq,j}(1-\xi_{\eq,j})}{w_j}
\quad(\text{식~\eqref{eq:lco-eqpeak} 의 피가수})\;,}
\label{eq:lco-msmrpeak}
\end{equation}
($|f|=1$; $\theta(1-\theta)=\xi(1-\xi)$ --- 여집합 불변).

\subsection{두 변경과 전자항 plug-in 사슬}\label{sec:lco-code-plugin}
따라서 Ch1 의 곡선 골격(평형 진행률 logistic$\cdot$평형 중심$\cdot$합산식~\eqref{eq:sum})은 구조 변경 없이
LCO 에 적용되며, 바뀌는 것은 단 둘이다:
\begin{enumerate}[label=(\roman*),leftmargin=2.2em,itemsep=1pt]
\item \textbf{전이 파라미터 교체} --- 흑연 초기값(표~\ref{tab:staging}) $\to$ LCO 초기값(표~\ref{tab:lco-staging}):
$(\Delta H_\rxn,\Delta S_\rxn,Q,\Omega,\Delta H_a,\dots)$ 값을 양극 영역으로.
\item \textbf{전자 엔트로피 항 plug-in} --- T1 전이의 $\Delta S_{\rxn}$ 평가에 몰당 전자항을 더하는 한 줄이며,
그 forward 사슬은 조성 좌표에서 dQ/dV 까지 아래와 같이 닫힌다.
\textbf{(a) 출발 --- 좌표 매핑의 닫힌 꼴.} 게이트(식~\eqref{eq:ggate})는 조성 $x$ 의 함수인데 곡선 계산은
전위점 $V$ 에서 직접 진행되므로, T1 창의 $x$ 범위(표~\ref{tab:lco-staging}: $x_{\mathrm{hi},1}{\approx}0.94$,
$x_{\mathrm{lo},1}{\approx}0.75$)를 T1 진행률로 선형 보간해 잇는다:
\begin{equation}
x\big(\xi_{\eq,1}(V)\big)=x_{\mathrm{hi},1}-\big(x_{\mathrm{hi},1}-x_{\mathrm{lo},1}\big)\,\xi_{\eq,1}(V)
\label{eq:lco-xmap}
\end{equation}
($\xi_{\eq,1}$ 은 탈리튬화-형 진행률(방향 규약은 \S\ref{sec:lco-direction}) --- $\xi{=}0\leftrightarrow x_{\mathrm{hi},1}$, $\xi{=}1\leftrightarrow x_{\mathrm{lo},1}$, $x_\mathrm{MIT}{\approx}0.85$ 가 창 내부; 표 두 끝점만 쓰는 단순 선형보간(tier C)이며 함수형 확정은 round-trip 피팅 몫).
\textbf{(b) 연산 --- 게이트 대입.} 식~\eqref{eq:lco-xmap} 를 게이트 인자에 넣고 몰당 나눗셈 형
(식~\eqref{eq:dSegate}$\cdot$\eqref{eq:gunit}$\cdot$\eqref{eq:dSemolar})으로 적으면 전위점 $V$ 의
전자항이 닫힌다:
\begin{equation}
\Delta S_{e,1}^{\,\mathrm{mol}}(V,T)
=-\frac{\pi^2}{3}\,R\,\frac{k_BT}{e_V}\,\frac{g_{\max}}{\Delta x_\mathrm{MIT}}\,
\sigma\!\big(z_e(V)\big)\big[1-\sigma\!\big(z_e(V)\big)\big],
\qquad z_e(V)=\frac{x(\xi_{\eq,1}(V))-x_\mathrm{MIT}}{\Delta x_\mathrm{MIT}}.
\label{eq:lco-SeV}
\end{equation}
\textbf{(c) 중간식 --- 슬롯 합과 온도 적분.} 이를 분해식~\eqref{eq:lco-decomp} 의 셋째 슬롯에 넣고
(★단위 주의 --- forward 슬롯 $\Delta S_{\rxn,j}$ 는 J/(mol\,K) 라, 자리당 식~\eqref{eq:dSe} 가 아니라
$N_A$ 를 곱한 몰당 식~\eqref{eq:dSemolar} 를 넣어야 한다; $N_A$ 배 누락 시 전자항이 $\sim10^{23}$ 배
과소), 온도 계수 관계~\eqref{eq:lco-dUdT} 를 기준온도에서 적분하면 T1 중심의 온도 이동이 닫힌다:
\begin{equation}
\begin{aligned}
\Delta S_{\rxn,1}^\mathrm{cat}(V,T)&=\Delta S_1^\mathrm{config}+\Delta S_1^\mathrm{vib}
+\Delta S_{e,1}^{\,\mathrm{mol}}(V,T),\\
U_1(V,T)&=U_1(T_\mathrm{ref})+\frac1F\int_{T_\mathrm{ref}}^{T}\Delta S_{\rxn,1}^\mathrm{cat}(V,T')\,\dd T'.
\end{aligned}
\label{eq:lco-U1V}
\end{equation}
전자항이 $\propto T'$ 라 이 적분의 닫힌형이 식~\eqref{eq:U1T2} 의 $T^2$ 곡률이고, \emph{현행 모델}은
$\Delta S_e$ 를 $T_\mathrm{ref}$ 에서 동결한 상수 오프셋(단일-기준 근사 --- 조성도 $x{=}x_\mathrm{center}$ 로
동결해 $V$-무관 상수)
으로 넣으므로 $U_1(T)\approx U_1(T_\mathrm{ref})+[\Delta S_{\rxn,1}^\mathrm{cat}(x_\mathrm{center},T_\mathrm{ref})/F]
(T-T_\mathrm{ref})$ 의 \emph{선형} $U(T)$ 가 된다(동결 근사의 결과 --- 다온도 $T^2$ 곡률은 round-trip 피팅 단계 과제). 평가 순서 주의 --- 식~\eqref{eq:lco-xmap} 이
$\xi_{\eq,1}$ 을, $\xi_{\eq,1}$ 이 $U_1$ 을 참조하는 고정점 구조이나, 동결 근사는 순환이 없고, 정밀형도
전자항 제외 중심으로 $\xi_{\eq,1}^{(0)}$ 을 먼저 평가해 1회 갱신하는 것을 초기 근사로 채택한다(되먹임
이동 $|\Delta S_e^{\,\mathrm{mol}}|/F\times|T-T_\mathrm{ref}|\approx0.47\,\mathrm{mV/K}\times30\,\mathrm K
\approx14$ mV $\lesssim w_1$ --- 폭 이하 스케일이라 보정이 작다는 상한이며, 수렴 여부 자체는 round-trip
피팅 단계에서 수치 확인한다). \textbf{(d) 박스 --- plug-in forward 사슬.}
\begin{equation}
\boxed{\;\begin{aligned}
x\big(\xi_{\eq,1}(V)\big)\ &\xrightarrow{\eqref{eq:lco-SeV}}\ \Delta S_{e,1}^{\,\mathrm{mol}}(V,T)
\ \xrightarrow{\eqref{eq:lco-U1V}}\ U_1(T)\ \xrightarrow{\eqref{eq:lco-Ubranch}}\ U_1^{\,d,\mathrm{cat}}\\[2pt]
&\xrightarrow{\eqref{eq:xieq}}\ \xi_{\eq,1}\ \xrightarrow{\eqref{eq:lco-eqpeak}}\
\Big(\frac{\dd Q}{\dd V}\Big)_{\!1}=Q_1\frac{\xi_{\eq,1}(1-\xi_{\eq,1})}{w_1}
\end{aligned}\;}
\label{eq:lco-plugin}
\end{equation}
$g(E_F,x)$ 는 식~\eqref{eq:ggate} 의 게이트로, 초기값 3개$(g_{\max},x_\mathrm{MIT},\Delta x_\mathrm{MIT})$
를 피팅 인자로 노출.
\end{enumerate}
$\partial U_j/\partial T\leftarrow\Delta S_{\rxn,j}^\mathrm{cat}/F$ 의 경로(식~\eqref{eq:lco-dUdT})는 흑연과 동일 ---
곧 ``파라미터 +1 항'' 외에 구조 변경이 없다는 것이 ``흑연 forward 골격이 LCO 양극까지 한 틀로 닫힌다''는
주장의 실증적 근거이자, 본 챕터가 두 전극을 한 프레임으로 통합한 까닭이다.

% 자산: [B2-112] [B2-113] [B2-114] [B2-115] [B2-116] [B2-117] [B2-118] [B2-119] [B2-120] [B2-121] [B2-122] [B2-123] [B2-124] [B2-125] [B2-126]
